\documentclass[11pt]{article}
\usepackage{amsmath,amssymb,amsthm,hyperref}
\newtheorem{theorem}{Theorem}\newtheorem{proposition}[theorem]{Proposition}
\newcommand{\E}{\mathcal E}\newcommand{\RR}{\overline{\mathbb R}_{\geq0}}
\title{Countable clock trace and causal execution: a native assembly theorem and its missing generator}
\date{16 November 2026}
\begin{document}\maketitle
\paragraph{Audit and scope.} The written construction has passed independent review at the declared scope; no proof-assistant certification is claimed. The construction relates instantaneous diagonal trace to causal branching execution through finite first-hit layers and countable assembly. It does not equate general trace with a least fixed point. A typed restriction operation is an explicit additional generator. A prefix-flow criterion characterizes which canonical terminal-layer weights admit the required causal origin; an arbitrary diagonal presentation need not do so. All terminal raw zero witnesses, actual augmented coefficients and the full unused terminating-tree grade are retained. No completion of the general v0.2 contract is claimed.

\section{Sources, coefficient doctrines and changes of ambient category}
The pinned inputs are the local retained copies of \emph{Native partial-trace domains, staged elimination, and shared coupling} (PR86), \emph{Variable-arity feedback through scope-atomic causal frontier strategies}, and the governing v0.2 extension contract. Their exact hashes are recorded separately. The first uses the native finite-measure writer trace; the second uses conditional frontier couplings and terminating trees at the full grade $R_UC$. This manuscript proves an assembly relationship, not a new unrestricted polynomial crossing.

Take countable external input, output and loop-state sets $A,B,U$. A primitive raw row is countable and normalized, with actual augmented context coefficients in $\E=\mathrm{Vect}_k/k$. Finite-dimensional individual coefficients from the source embed fully faithfully; finite products of them remain finite-dimensional. Countable vector-space direct sums exist for row aggregation in Vect, whereas their existence in FinVect must not be presumed. The temporal bind used here is the source's directly constructed product-coefficient bind, not an unsupported distributive rig structure on $\E$: at a paired witness its coefficient is $E\times F$, with weight the product. Actual maps are products of the supplied arrows. The original tensor-coefficient independent semantics maps to this doctrine through its augmentation-induced tensor-to-product comparison, generally not invertibly.

The independent writer category $\mathcal D$ is symmetric monoidal. Its products below are finite tensor products, so no infinite tensor product or countable coproduct in $\mathcal D$ is assumed. A global grade is a full small family $(O,(D_o))$, with all unused labels retained. An arity-annotated tree grade is encoded as that same writer family and an ordered output list in $B^*$ satisfying $|\mathbf b|=\operatorname{arity}(o)$. This is a faithful typed encoding of $F_{R_UC}(B)$; it does not discard the arity or any collection map. The set $B^*$ is countable. Native grade maps in this subinterface preserve arities.

\paragraph{Lossless polynomial-body encoding at the stated cells.}
For a fixed full collection $O=(O_n)$, the map
\[
 \coprod_n O_n\times B^n\ \longrightarrow\ O\times B^*,
 \qquad(o,(b_1,\ldots,b_n))\longmapsto(o,[b_1,\ldots,b_n])
\]
is bijective onto the length--arity predicate. Its inverse reads $n$ from $o$ and the ordered list. A native body square consists of an ordered arity-preserving collection map $\alpha$, an endpoint map $b:B\to B'$, and its actual coefficient arrows. Under this bijection its equation is exactly the writer body equation with list map $b^*$. Conversely such a square recovers every component of the polynomial square. Thus the encoding is fully faithful on these specified cells and reflects their equality. Arbitrary maps $B^*\to(B')^*$, arbitrary arity-changing writer maps, and arbitrary writer objects outside the predicate are not part of the claimed equivalence. The clock trace preserves the predicate because it changes neither tree shape nor leaf list.

\paragraph{Countable coefficient folds are actual, unweighted maps.}
A coefficient row can be aggregated as $\bigoplus_w E_w$ with augmentation to $k$ given on each summand by its augmentation. A raw map $r$ and augmented arrows $\phi_w:E_w\to E'_{r(w)}$ induce the actual arrow from that direct sum to $\bigoplus_{w'}E'_{w'}$ determined by those components. Every input vector has finite support, so even a countably infinite fibre requires no infinite vector sum. Composition is equality on summands; augmentation preservation follows on each summand. Scalar weights, including infinity, form a separate component and are never used as vector-space scalar multipliers. Dependent bind uses the displayed individual product coefficients and the product of actual maps, so no countable distributivity of $\E$-products over these direct sums is required. Independent tensor coefficients, if separately used, distribute over direct sums in Vect, but they are not identified with these products.

For the ambient countable trace allow weights in $\RR=[0,\infty]$, with $0\cdot\infty=0$, and countable raw witness sets. Every countable sum is the supremum of finite subsums. A native arrow is a raw function satisfying exact fibre pushforward, together with actual coefficient arrows and the full global writer map and body square. This is an explicit enlargement of scalar effects. Finite-total but unbounded row masses alone do not form a composition-closed alternative: a row with weights $2^{-n}$, $n\geq1$, followed at outcome $n$ by a singleton of mass $2^n$ has infinite composite mass.

\section{Countable ambient trace and summability domains}
Composition uses dependent raw pairs, product weights, product context coefficients, and tensor writer grades. Tensor uses independent raw pairs, product weights, categorical-product context coefficients, and tensor writer grades. Countability is preserved by these constructions. Nonnegative countable Fubini and finite-product coherence prove associativity, units and interchange, including infinite weights. Fubini here follows directly by taking suprema over finite subsets of the countable index sets; multiplication by a fixed extended scalar preserves these sums with the stated zero convention.

For $F:X\times V\to Y\times V$ with countable $V$, define
\[
 W^{\operatorname{tr}}_x=\coprod_{v\in V}\{w\in W_{x,v}:v^+(w)=v\},\qquad
 \mu^{\operatorname{tr}}_x(v,w)=\mu_{x,v}(w).                 \tag{1}
\]
Coefficients, external bodies and the entire global grade are unchanged. This is an operation on the actual native presentations, not a sum of their carrier interpretations.

\begin{theorem}[Countable matching trace]
Formula (1) is the native compact trace in the extended-weight countable-endpoint writer bicategory. Its restriction to subprobability processes is a partial trace with domain
\[
 z_F(x):=\sum_{v,w:v^+(w)=v}\mu_{x,v}(w)\leq1
 \quad\hbox{for every }x.                                \tag{2}
\]
Its restriction to normalized processes has domain $z_F(x)=1$. The partial-trace axioms have directed naturality and strength, bidirectional dinaturality, and vanishing-II conditional on the inner trace already being admitted, exactly as in PR86. Fixed-loop native cells preserve the quantities $z_F$ and induce actual traced cells. General loop-changing native squares are not claimed to be preserved.
\end{theorem}
\begin{proof}
The cup on $V$ has raw witness $v$ of weight one and output $(v,v)$; its total mass may be infinite, allowed in this ambient category. The cap has a singleton weight-one witness on equal inputs and no witness otherwise. The snake composites have exactly one compatible raw witness of weight one for each input. Actual unitors give their coefficient and grade isomorphisms. The cup--cap trace has exactly the witnesses in (1), with the actual coefficient and grade unitors removing inserted units. Associativity, dinaturality, superposing and nested trace compare the same countable tuples and matching equations; Fubini permits their regrouping and symmetric monoidal coherence supplies the actual arrows. No cancellation of infinity is used.

A fixed-loop body square reflects as well as preserves the matching predicate. Exact fibre pushforward can therefore be restricted to matching witnesses, retaining all zero witnesses and all coefficient maps, and proves preservation of $z_F$. Subprobability composition and tensor are closed by Fubini. If $G$ and $K$ are subprobability pre- and post-processes, the closed mass of $(K\times1)F(G\times1)$ is at most the $G$-weighted average of $z_F$; hence (2) implies its admission. Tensoring with a subprobability process likewise preserves the bound. For normalized $G,K$, the analogous equality proves the normalized assertion. Dinaturality's tuple bijection gives equal total masses in both presentations. Provided the inner trace is in the selected subcategory, nested and joint trace have equal rows and total masses, so their admission is equivalent. Vanishing-I and yanking have their prescribed admitted values. These arguments restrict the actual native coherence maps, not just numerical observations.
\end{proof}

\section{Finite causal generation, with restriction made explicit}
Fix the source causal-frontier interface. A history $h$ retains the full addressed raw tree and current ordered frontier $I(h)$. At every nonterminal history, including zero histories, a strategy supplies a normalized joint law $q_h$ on the finite product of primitive raw rows, with prescribed primitive marginals. Its successor $h\cdot\mathbf w$ is determined by the primitive bodies. Each transition adds one complete frontier round. A terminal history has empty frontier and never advances. All finite histories from one external root are countable. Put
\[
 m(a)=1,\qquad m(h\cdot\mathbf w)=m(h)q_h(\mathbf w),\qquad
 E_h=\prod_{v\in\mathrm{nodes}(h)}E_{x_v,w_v}.              \tag{3}
\]
These are finite products for each history, even when the set of histories is infinite.

For $n\geq1$, form the finite-horizon prefix object by applying the supplied kernels for at most $n$ rounds and making earlier terminal histories absorbing for this bookkeeping only. It retains every syntactic raw prefix, including zero ones, and its prefix-tree data and coefficient. Its total mass is one. Then take the typed stratum of histories whose \emph{first} terminal round is exactly $n$; denote it $T_{a,n}$. The layer $L_n(a)$ has this raw set, weights (3), coefficients $E_t$, and its actual tree body at the full global grade $R_UC$.

\paragraph{Restriction is not a hidden native arrow.}
The prefix object and its rejected stratum remain distinct stored constructions. The passage to $T_{a,n}$ is a typed raw-domain restriction, not an exact-pushforward arrow from the entire prefix object. Such an arrow would in general falsely preserve total mass. For example, a fair primitive row with one terminal and one continuing raw witness has total mass one, whereas its first-round terminal stratum has mass $1/2$; summing the exact-pushforward equations would require $1=1/2$. It is a separately supplied partial-observation operation on typed rows. A convenient precise interface takes a row together with a specified subset of its raw set and returns the restricted row; maps are allowed only when the selected subset is the inverse image of the target selected subset. Restriction then is functorial because the native fibre equation restricts to exactly the selected fibres. The target tree body is required only on the selected subset; no nonexistent terminal body is assigned to an active prefix. Restriction retains every selected zero-weight point. Thus no discarded active prefix is misdescribed as a zero terminal witness.

\begin{theorem}[Causally generated clock trace]
From finite unrolling, the stated typed restriction, and countable matching trace, define the clock process
\[
 K_\sigma:A\times\mathbb N_{>0}\longrightarrow B^*\times\mathbb N_{>0}
\]
whose row at $(a,n)$ is $L_n(a)$ and whose body returns $(\mathbf b(t),n)$, with the same full writer grade $R_UC$. Each row is subprobabilistic. Its countable trace is admitted and is natively isomorphic, at identity full grade, to the source terminating execution evaluator $f^{\dagger,\sigma}$. This compiler uses only finite products and finite-horizon kernel assembly to define each clock row; it never queries the already assembled terminating law.
\end{theorem}
\begin{proof}
Induction on the finite horizon proves that terminal-by-$n$ mass plus still-active mass is one. Each transition distributes a prefix's mass by a normalized kernel; countable summation commutes with this distribution. The first-terminal strata are disjoint. Therefore $\sum_{j\leq n}\sum_{t\in T_{a,j}}m(t)\leq1$, and taking the supremum proves (2) for $K_\sigma$. Every finite terminal execution has one and only one first-terminal height. The map $(n,t)\mapsto t$ is consequently a bijection from the traced raw set to all finite terminal histories, including every zero-weight terminal history. Formula (3) is exactly the product of the successive causal probabilities in the evaluator. Its coefficient product has exactly the same node factors; the canonical product associator gives an actual isomorphism. The ordered leaves and full tree body agree. The global map is identity on \emph{all} $R_UC$, not its realized image. These maps and their inverses are therefore actual native arrows.
\end{proof}

\section{Native transport and stopping-cut coherence}
A primitive native comparison preserving ordered arity and the terminal/continuing tag induces a history map $H$ preserving exact height and the terminal predicate. Assume the source's local conditional compatibility
\[
 (r_h)_*q_h=q'_{Hh}                                       \tag{4}
\]
at every source history. Histories identified by $H$ must have that same target conditional kernel; no quotient by almost-sure equality is taken.

\begin{theorem}[Assembly naturality]
The compiler and its clock-trace isomorphism are natural for these comparisons. They retain the entire grade map $R_u\alpha$, every actual node coefficient arrow, and unused shapes. They commute with bounded stopping-frontier cut/assembly at the full marked-cut grade; forgetting the cut is the specified directed full-grade map, not an isomorphism. Repeated admitted cuts commute by the same native coherence.
\end{theorem}
\begin{proof}
Induction on horizon proves exact pushforward of prefix mass: for a target successor, first sum over the source frontier draws above its draw, using (4), and then over source prefixes, using the inductive equation. Because $H$ preserves height and the terminal tag, the exact-$n$ terminal stratum is the inverse image of that stratum. The restriction operation therefore transports the native map rather than discarding an uncontrolled fibre. Taking the countable disjoint union over $n$ preserves each exact fibre equation. The loop clock map is identity, so countable trace gives the same arrow. On coefficients both routes have the same projection to every target node, followed by its supplied node map; finite-product uniqueness proves equality. On grades both are the same full $R_u\alpha$.

For a bounded stopping barrier, each terminal history splits uniquely as a cut prefix $k$ and a joint suffix forest $s$. Its first-terminal height is $d(k)+\ell(s)$; the suffix length is zero if $k$ is already terminal. Its weight is the prefix weight times the conditional suffix weight, and its coefficient is canonically $E_k\times E_s$. Thus tracing the clock layers of the whole execution and summing the dependent prefix/suffix layers give the same raw set via $(n,t)\leftrightarrow(k,\ell(s),s)$ with the stated relation on lengths. There are no extra idle rounds. Countable Fubini proves equality of the resulting weights. Retaining the cut marker gives exactly the source's full $\operatorname{Cut}_UC$ grade, and both constructions use the actual graft/forget-cut map to $R_UC$ on all marked shapes. A compatible native comparison preserves the stopping decision and commutes with this bijection. With multiple cuts all routes partition the same finite list of rounds and node factors; product uniqueness and the identical full marked-grade maps prove coherence. This is not a cut through an atomic frontier sampling scope.
\end{proof}

\section{Which terminal layers have a causal origin?}
The preceding theorem is stronger than the universal tautology that any countable weighted family is its own diagonal trace. An independently checkable extension criterion distinguishes causal origin. Fix the primitive bodies and rows and hence the entire syntactic history tree. A proposed terminal-layer family assigns nonnegative weights $b_t$ to \emph{all} its terminal histories, retains the canonical node-product coefficients and tree bodies, and retains the full grade $R_UC$. No positive-support quotient is allowed.

\begin{theorem}[Prefix-flow characterization]
Such a family is the clock trace generated by some causal-frontier strategy if and only if there exist nonnegative numbers $m_h$ on all finite histories, with root mass one, terminal masses $m_t=b_t$, and for every nonterminal $h$ the equations
\[
 \sum_{\mathbf w}m_{h\cdot\mathbf w}=m_h,\qquad
 \sum_{\mathbf w:w_i=w}m_{h\cdot\mathbf w}=m_h p_{x_i(h)}(w)
 \quad(i\in I(h),\ w\in W_{x_i(h)}).                     \tag{5}
\]
These are linear nonnegative countable-flow constraints involving only primitive marginals and the syntactic prefix relation. Any such flow has $0\leq m_h\leq1$ and terminal total at most one. Over a fixed feasible flow, the complete strategies realizing it are exactly the forced conditional kernels $m_{h\cdot\mathbf w}/m_h$ at positive histories and arbitrary points of the prescribed marginal-coupling set (a polytope when the rows are finite) at every zero history.
\end{theorem}
\begin{proof}
A strategy supplies (3). Summing the normalized kernel gives the first equation, and its specified marginals give the second. Conversely, nonnegativity and conservation imply recursively that every child mass is at most its parent and all masses are finite. At a positive history define $q_h(\mathbf w)=m_{h\cdot\mathbf w}/m_h$; (5) proves its normalization and all required marginals. At a zero history, conservation forces every child mass to be zero. Choose there any normalized coupling of the primitive rows, for instance their finite independent product. This is available even for countable raw rows. Induction then recovers precisely every prescribed $m_h$, including zero histories and terminal masses. The characterization of the fibre follows because division is forced exactly at positive histories and imposes no constraint at zero histories. The usual finite-height conservation proof gives the terminal-mass bound. Coefficients and body maps need no reconstruction: their retained canonical finite-node definitions are already the required ones. There is no claim of functorial arbitrary selection of zero-history couplings.
\end{proof}

This is a characterization of the image at the stated canonical native object interface, not a universal property of all feedback or a reconstruction from terminal weights alone. Feasible flow can fail: if a primitive root terminates with two distinct raw witnesses of probabilities $1/2,1/2$, a proposed normalized terminal family of masses $1,0$ has a perfectly valid singleton-clock diagonal trace but violates (5) at the root. Feasible flow is also not generally unique, because terminal observations do not determine the allocation of nontermination among active prefixes. Even a fixed feasible flow loses all zero-history strategy controls unless they are separately retained. The flow criterion therefore explains precisely why arbitrary path-index dilations are insufficient evidence of causal generation.

\section{Grade and mechanism obstructions}
\begin{proposition}[No generic normalized clock filler at unchanged full grade]
There is no generic replacement of the typed first-hit restriction by a normalized clock row at the same full grade and endpoint, valid for every primitive causal system.
\end{proposition}
\begin{proof}
Take $B=\varnothing$, one loop state, and only a unary primitive shape, always returning the loop state. Every full finite nonempty tree in $R_UC$ then has one external leaf; there are no nullary shapes. Hence $F_{R_UC}(B)=\varnothing$. No normalized row can have a body in this empty set, for any clock input. The terminating evaluator nevertheless exists: its raw set is empty and its full nonempty grade is still $R_UC$. The killed first-hit layers are empty subprobability rows and compile correctly. Adding a cemetery shape would permit a normalized rejected outcome, but native trace preserves that larger entire grade; it cannot silently delete the cemetery shape to recover the claimed full-grade object. A separately supplied appropriate grade map and body would be additional data, not a generic solution; here no arity-preserving map can send a new nullary cemetery shape into the original unary-only collection.
\end{proof}

Countable summability is genuine: a primitive that terminates with probability $p\in(0,1)$ and otherwise repeats has terminal layer mass $p(1-p)^{n-1}$. Every finite union of layers misses positive mass and raw histories; countable clock trace recovers total mass one while retaining all finite-history coefficients. No infinite coefficient tensor is needed. Conversely, a primitive that always repeats has zero terminating trace, although its finite-prefix system has mass one at every horizon. The construction does not assign a particular divergent trajectory or equate missing mass with a designated refusal witness.

\paragraph{Recurrent prefix information and its coefficient boundary.}
Before restriction, the absorbed finite-horizon prefix laws form a coherent probability system under raw truncation. For any target prefix, summing all extensions to the next horizon gives its mass by normalization of the causal kernel, including zero-prefix fibres. Truncation has an actual context-coefficient arrow from the longer node product to the shorter one, given by its retained-node projections. Comparisons satisfying (4) commute with these truncations on raw histories and context coefficients. This retains coherent recurrent information even when the terminating trace is zero.

It does not yet supply a pro-object of the original full polynomial native processes: truncating a partial tree can change leaf arity, and its full content coefficient needs deletion of descendant tensor factors. A merely symmetric monoidal $\mathcal D$ need not supply arrows $D\to I$ for that deletion. Additional chosen content augmentations and a typed partial-tree truncation interface are needed before claiming that stronger graded pro-native object. Likewise a scalar infinite-path probability measure would not by itself construct an accumulated infinite coefficient or full native infinite-execution object. This precise residual is not repaired by the countable trace theorem, whose retained histories are finite.

The independently supplied operations are now visible. Instantaneous trace matches and sums an already typed family. Causal kernels construct finite prefixes. Typed restriction selects first-hit terminal outputs. Countable assembly combines their layers, with the flow criterion deciding causal realizability. Removing any of these distinctions would turn the theorem into an arbitrary diagonal encoding or an incorrect claim that trace chooses a least fixed point. A separate coproduct-based path-sum iteration doctrine may implement the restriction differently, but it must still account for full grades, partial bodies and native zero witnesses. No equivalence with such a doctrine, unrestricted shared polynomial substitution, or probability law on infinite executions is asserted here.
\end{document}
